\documentclass[a4paper,12pt]{article}
\usepackage[utf8]{inputenc}
\usepackage[russian]{babel}
\usepackage{amsmath, amssymb}

\title{Лекция 6 \\ Циклические коды}
\author{}
\date{}

\begin{document}

\maketitle

\section{Циклические коды}

Циклические коды --- это подкласс линейных блоковых кодов, обладающий
дополнительной алгебраической структурой: циклический сдвиг любого
кодового слова также является кодовым словом.

Формально, если
\[
(c_0, c_1, \dots, c_{n-1}) \in C,
\]
то
\[
(c_{n-1}, c_0, \dots, c_{n-2}) \in C.
\]

Кодовые слова удобно представлять в виде полиномов степени не выше $n-1$
над полем $\mathbb{F}_q$ (чаще всего $\mathbb{F}_2$), рассматривая арифметику
по модулю $x^n - 1$.

\subsection{Свойства}

\begin{itemize}
    \item Циклические коды являются идеалами в кольце $\mathbb{F}_q[x]/(x^n - 1)$.
    \item Каждый код задаётся порождающим полиномом $g(x)$.
    \item $g(x)$ делит $x^n - 1$.
\end{itemize}

\subsection{Примеры}

\begin{itemize}
    \item коды Хэмминга,
    \item БЧХ-коды,
    \item коды Рида--Соломона.
\end{itemize}

Для RS-кодов:
\[
n = q - 1, \quad d = n - k + 1
\]
(достигается граница Синглтона).

\section{Теорема об умножении на полином}

Пусть имеется циклический $(n,k)$-код и $g(x)$ --- кодовое слово.
Тогда для любого полинома $m(x)$:
\[
m(x)g(x) \mod (x^n - 1)
\]
также является кодовым словом.

\subsection*{Доказательство}

Циклический сдвиг соответствует умножению на $x$ по модулю $x^n - 1$.
По линейности любая линейная комбинация таких сдвигов также принадлежит коду,
что и даёт требуемый результат.

\subsection*{Пример}

\[
G =
\begin{pmatrix}
1 & 1 & 0 \\
0 & 1 & 1
\end{pmatrix}
\]

Строки соответствуют полиномам:
\[
1 + x, \quad x + x^2
\]

Линейная комбинация:
\[
a(1,1,0) + b(0,1,1) = (a, a+b, b)
\]

Полиномиально:
\[
a + (a+b)x + bx^2
\]

\section{Теорема о существовании порождающего полинома}

В любом циклическом $(n,k)$-коде существует единственный приведённый
полином $g(x)$ минимальной степени.

\subsection*{Доказательство}

Выберем ненулевой полином минимальной степени. Если бы существовал
другой полином той же степени, их разность имела бы меньшую степень,
что противоречит минимальности.

\section{Теорема о делимости кодовых слов}

Любое кодовое слово $c(x)$ делится на $g(x)$:
\[
c(x) = m(x)g(x)
\]

\subsection*{Пример}

\[
g(x) = 1 + x + x^3, \quad n = 7
\]

\[
c(x) = 1 + x^2 + x^3 + x^4
\]

\[
c(x) = g(x)(1 + x + x^2)
\]

\section{Теорема о делимости $x^n - 1$}

Порождающий полином делит $x^n - 1$:
\[
x^n - 1 = g(x)h(x)
\]

\subsection*{Идея доказательства}

Если остаток при делении был бы ненулевой, он дал бы кодовое слово
меньшей степени, что невозможно.

\section{Теорема о степени порождающего полинома}

Для циклического $(n,k)$-кода:
\[
\deg g(x) = n - k
\]

\subsection*{Доказательство}

Код состоит из:
\[
c(x) = m(x)g(x), \quad \deg m(x) < k
\]

Отображение $m(x) \mapsto m(x)g(x)$ инъективно, следовательно
размерность пространства равна $k$, откуда:
\[
\deg g(x) = n - k
\]

\end{document}